% Part:sets-functions-relations
% Chapter: sets
% Section: schroder-bernstein

\documentclass[../../../include/open-logic-section]{subfiles}

\begin{document}

\olfileid[tr]{sfr}{siz}{sb}

\olsection{Büyüklük Kavramı ve Schr\"oder--Bernstein}

\begin{explain}
Şöyle sezgisel bir düşünce vardır: $A$, $B$'den daha büyük değilse ve $B$ de
$A$'dan daha büyük değilse $A$ ile $B$ eş güçlüdür. Açıkçası bu düşünce
\emph{yanlış} olsaydı, tanımladığımız eş güçlülük kavramının kümeler
arasındaki ``büyüklük'' karşılaştırmalarıyla bir ilgisi olduğunu savunmakta
çok zorlanırdık!{} Neyse ki sezgisel düşünce doğrudur. Bunu
Schr\"oder--Bernstein Teoremi temellendirir.
\end{explain}

\begin{thm}[Schr\"oder--Bernstein]
	\ollabel{thm:schroder-bernstein}
	$\cardle{A}{B}$ ve $\cardle{B}{A}$ ise $\cardeq{A}{B}$.
\end{thm}

\begin{explain}
Başka bir deyişle $A$'dan $B$'ye ve $B$'den $A$'ya birer !!{injection}
varsa $A$'dan $B$'ye bir !!{bijection} vardır.

Ne var ki bu sonucu ispatlamak gerçekten oldukça \emph{zordur}. Gerçekten
sonucu Cantor ifade etmiş, başkaları ispatlamıştır.\footnote{Tarihçesi için
örneğin \citet[pp.~165--6]{Potter2004}'e bakın.}
\oliflabeldef{sfr:cardinals:card-sb:sec}{Schr\"oder--Bernstein'ı ancak
\olref[sfr][cardinals][card-sb]{sec}'te \emph{ispatlayabilecek} durumda
olacağız.}{}%
Şimdilik ona güvenebilirsiniz (ve güvenmeniz gerekir).

Neyse ki Schr\"oder--Bernstein \emph{doğrudur} ve tanımladığımız
$\cardeq{A}{B}$ ile $\cardle{A}{B}$ bağıntılarının ``büyüklük''le ilgili
olduğu düşüncemizi doğrular. Dahası Schr\"oder--Bernstein çok
\emph{yararlıdır}. Eş güçlü iki küme arasında bir !!{bijection} bulmak zor
olabilir. Schr\"oder--Bernstein Teoremi karşılaştırmayı iki yöne bölmemizi;
yalnızca birinciden ikinciye ve tersine birer !!{injection} düşünmemizi sağlar.
\end{explain}
\end{document}
